% =============================================================================
% blocks.tex —— 三色圆角 Block（普通 / 警示 / 示例）
% -----------------------------------------------------------------------------
% 目标：给 Beamer 原生圆角盒加上内边距，并让块内的列表符号与 \alert
%       跟随该块自己的颜色（而不是全局主题蓝）。
% 做法：不引入任何新宏包，只用 beamerboxesrounded + minipage 内缩。
% 依赖：colors.tex 必须先加载（用到 block title/body 三套配色）。
% 调参入口：\InkBlockPadding 控制水平内边距；两个 \vspace 控制竖直内边距。
% =============================================================================

% 标题栏与正文的左右内边距
\newlength{\InkBlockPadding}
\setlength{\InkBlockPadding}{0.5em}

% 内层 minipage 的可用宽度 = 行宽 - 左右内边距（只需算一次，避免重复算式）
\newlength{\InkBlockInnerWidth}

% 把 block 内的「局部结构色」和「强调色」重绑到当前块的标题色，
% 于是块内的 itemize/enumerate 符号和 \alert 会自动变成该块的颜色。
% 参数：空 / " alerted" / " example"（注意前导空格，用于拼接模板名）
\newcommand{\InkBlockColors}[1]{%
  \setlength{\parskip}{0.3em}%
  \setbeamercolor{local structure}{use={block title#1},fg=block title#1.bg}%
  \setbeamercolor{alerted text}{use={block title#1},fg=block title#1.bg}}

% 块开始：先排一个自带内边距的标题栏，再开启正文 minipage。
% shadow=false 与规范所述「无阴影」一致；width=\linewidth 让块自适应单栏或分栏。
\newcommand{\InkBlockBegin}[1]{%
  \par\vskip\medskipamount
  % 原生圆角盒在左右各向外延伸 4bp；内缩后让外框与正文边界对齐。
  \noindent\hspace*{4bp}%
  \begin{minipage}{\dimexpr\linewidth-8bp\relax}
  \setlength{\InkBlockInnerWidth}{\dimexpr\linewidth-\InkBlockPadding-\InkBlockPadding\relax}%
  \begin{beamerboxesrounded}[upper=block title#1,lower=block body#1,width=\linewidth,shadow=false]{%
    \noindent\hspace*{\InkBlockPadding}%
    \begin{minipage}{\InkBlockInnerWidth}
      \usebeamerfont{block title#1}\vspace*{0.2ex}%
      \strut\insertblocktitle\par\vspace*{0.2ex}%
    \end{minipage}}
    \usebeamerfont{block body#1}\InkBlockColors{#1}%
    \noindent\hspace*{\InkBlockPadding}%
    \begin{minipage}{\InkBlockInnerWidth}
      \vspace*{0.6ex}}

% 块结束：补上正文下边距并闭合两层盒子
\newcommand{\InkBlockEnd}{%
      \par\vspace*{0.6ex}%
    \end{minipage}\end{beamerboxesrounded}%
  \end{minipage}\par\vskip\smallskipamount}

% 把上面两个宏挂到 Beamer 的 six 个 Block 模板上
\setbeamertemplate{block begin}{\InkBlockBegin{}}
\setbeamertemplate{block end}{\InkBlockEnd}
\setbeamertemplate{block alerted begin}{\InkBlockBegin{ alerted}}
\setbeamertemplate{block alerted end}{\InkBlockEnd}
\setbeamertemplate{block example begin}{\InkBlockBegin{ example}}
\setbeamertemplate{block example end}{\InkBlockEnd}

% 块内文字：无衬线（黑体）粗体标题 + 常规正文；
% alerted / example 的字体继承普通块，只换配色。
\setbeamerfont{block title}{family=\sffamily,series=\bfseries}
\setbeamerfont{block body}{family=\sffamily,size=\normalsize,series=\mdseries,shape=\upshape}
\setbeamerfont{block title alerted}{parent=block title}
\setbeamerfont{block title example}{parent=block title}
\setbeamerfont{block body alerted}{parent=block body}
\setbeamerfont{block body example}{parent=block body}
